\documentclass[12pt]{article}
\usepackage[utf8]{inputenc}
\usepackage[T1]{fontenc}
\usepackage{hyperref}
\title{Yekora Language Resources for K-12 Curriculum Development}
\author{Compiled by Claude Sonnet 5.5 (claude-sonnet-5-5) for Andr\'as Kornai}
\date{}
\begin{document}
\maketitle

\section{Introduction}
Yekora (ISO 639-3 ykr, Glottocode yeko1240) is a Papuan language of the Binandere branch of Trans--New Guinea, spoken in Morobe Province, Papua New Guinea. Its closest relative is Zia. Wikipedia gives about 1,100 speakers (2000) and Global Recordings Network about 1,000. Joshua Project places the community south of Morobe Station along the coast and inland along the Mo River. Most speakers also use Zia and Suena. There is no vernacular Bible: church services use English or Tok Pisin scriptures. The search found no school materials.

\section{Primary Educational Resources}
The searches found no primers, textbooks or school programmes in Yekora.

\section{Linguistic and Reference Materials}
\begin{itemize}
\item Wikipedia, ``Yekora language''. \url{https://en.wikipedia.org/wiki/Yekora_language}
\item ASJP (Automated Similarity Judgment Program) wordlist entry for Yekora. \url{https://asjp.clld.org/languages/YEKORA}
\item Trans New Guinea database, Yekora entries. \url{https://transnewguinea.org/language/yekora?page=4}
\item MPI EVA numeral data for Yekora (from Darryl Wilson). \url{https://lingweb.eva.mpg.de/channumerals/Yekora.htm}
\end{itemize}
These sources are wordlists and numerals, not a grammar or dictionary.

\section{Religious and Community Resources}
\begin{itemize}
\item Joshua Project, Yekora in Papua New Guinea. \url{https://joshuaproject.net/people_groups/16016/PP}
\end{itemize}
A Yekora scripture translation has been proposed, based on the completed Zia or Suena New Testaments, but I found none.

\section{Digital and Audio Resources}
\begin{itemize}
\item Global Recordings Network, language page for ykr. I did not see its contents in the search results. \url{https://globalrecordings.net/en/language/ykr}
\end{itemize}

\section{Curriculum Development Context}
Yekora is stable but small, and speakers also use Zia and Suena. Documentation is limited to wordlists. A curriculum would need new orthography work, ideally with Zia as a comparison, along with community-created materials. English and Tok Pisin would carry the rest of the curriculum.

\end{document}